\documentclass[12pt]{article}

%%\usepackage[english]{babel}
\usepackage[T1]{fontenc}
\usepackage[polish]{babel}
\usepackage{amsmath}
\usepackage{amssymb}
\usepackage{url}


\begin{document}

\title{Econometrics Notes VI: Elementary Statistics - Part 3: Distributions - Exponential, \\ Gamma, $\chi^2$}
\author{Artur Wegrzyn}
\maketitle

\abstract{This note presents basic results concerning distributions that will show up frequently in later considerations: exponential, gamma and $\chi^2$. }

\tableofcontents


\section{Characteristic Function Uniquely Determines Probability Distribution}
In this note the fact that characteristic function uniquely determines probability distribution is used. 
For proof and comprehensive discussion of this theorem cf.:
\begin{itemize}
\item in English: \cite{feller_probability_vol2} - chapter XV: \textit{Characteristic Functions}
\item in Polish: \cite{jakubowski_sztencel} - chapter 9: \textit{Funkcje charakterystyczne}
\end{itemize}
In these notes, however, I will work with moment-generating functions instead of characteristic functions 
for the sake of clarity.

\section{Exponential Distribution}
Random variable $X$ follows exponential distribution with rate\footnote{Rate parameter is inverse of scale 
parameter. Recall that $s$ is a scale parameter if $F(x; s, \theta) = F(x/s; 1, \theta)$} parameter $\lambda$
if its pdf is:
\begin{equation*}
f_{X}(x) = \lambda e^{-\lambda x}
\end{equation*}
MGF of $X$ is:
\begin{equation*}
M_{X}(t) = E(e^{tX}) = 
\end{equation*}
The fact that $X$ is exponentially distributed with scale parameter lambda is denoted:
\begin{equation*}
X \sim Exp(1/\lambda)
\end{equation*}


\section{Gamma Distribution}

\subsection{Sum of $p$ Exponential Random Variables}
Let $X_1, X_2, \dots, Y_p$ be a $p$-element sample drawn indepentently from an exponential 
distribution with expected value $1 / \lambda$ i.e.:
\begin{equation*}
X_i \sim Exp(1/\lambda)
\end{equation*}

\subsection{•}


\section{Square of Standard Normal Distribution}
Suppose that $U$ follows standard normal distribution:
\begin{equation*}
U \sim N(0, 1)
\end{equation*}

Consider $Z = U^2$. Then:
\begin{equation*}
P(Z \leq z) = P(U^2 \leq z) = P(-\sqrt{z} < U \leq \sqrt{z})
\end{equation*}

Hence:
\begin{equation*}
f_{Z}(z) = \frac{dF_{Z}(z)}{dz} = \frac{d}{dz} P(-\sqrt{z} < U \leq \sqrt{z})
\end{equation*}

A little bit of algebra yields:
\begin{equation*}
f_{Z}(z) = 
\end{equation*}


\section{$\chi^2$ Distribution}
In this and subsequent section we consider a sample $X_1, X_2, \dots, X_n$ 
that is drawn independently from a normal 
distribution with mean $\mu$ and variance $\sigma^2$:
\begin{equation*}
X_i \overset{\text{i.i.d.}}{\sim} N(\mu, \sigma^2)
\end{equation*}

\bibliography{tsa_bibliography}
\bibliographystyle{ieeetr}


\end{document}
